% Part II, mechanising program calculation (files 05 and 06, about
% 20 min): the gap, and our answer (the survey table is in the backup)
\partframe{Part II}{Mechanising Program Calculation\\[6mm]
  {\normalsize Joint work with Marco Paviotti and Michael Vollmer\\[1mm]
  \textit{Program Optimisations via Hylomorphisms for Extraction of
  Executable Code}, ITP 2025}}

% \begin{frame}
%   \frametitle{The Gap in Mechanised Program Calculation}
%   \vfill
%   \begin{center}
%     \begin{minipage}{.8\textwidth}
%       \begin{sticky}
%         \textbf{Hylomorphisms in a proof assistant},
%         \begin{itemize}
%           \item<2-> for \textbf{any} data type, and recursion the
%             prover cannot see terminates,
%           \item<3-> with the laws as \textbf{regular equalities},
%           \item<4-> with \textbf{no axioms},
%           \item<5-> extracting \textbf{clean, runnable} code.
%         \end{itemize}
%       \end{sticky}
%     \end{minipage}
%   \end{center}
%   \vfill
%   \uncover<6->{Earlier mechanisations: lists only, or no divide and
%     conquer, or extra axioms, or no code.}\par\smallskip
%   \uncover<7->{\textbf{Our answer}: a Rocq library, joint work with
%     M. Paviotti and M. Vollmer (ITP 2025).}\par\vspace*{7mm}
%   \only<1-6>{\noslideref}%
%   \only<7->{\slideref{D. Castro-Perez, M. Paviotti, M. Vollmer. \textbf{Program Optimisations via Hylomorphisms for Extraction of Executable Code}. ITP 2025.}}
%   \speakernote{Generic: any polynomial shape, divide and conquer \\
%     Laws as regular equalities: rewrite with them, as on paper \\
%     No axioms: results hold in plain Rocq \\
%     Clean code: what you would write by hand, no casts \\
%     Each earlier work misses at least one (survey table: backup) \\
%     Our answer: the ITP 2025 library, with Marco Paviotti and Michael
%     Vollmer; the rest of this part}
% \end{frame}
